\section{Spatial-Spectral Audio Action Model (S2A2)}
	本節では，提案するマルチモーダル模倣学習フレームワークS2A2の全体構成を述べる．
	まず~\ref{sec:problem}で問題設定を定義し，~\ref{sec:tasks}で本研究が対象とする音環境マニピュレーションタスク群を導入する．
	続いて~\ref{sec:overview}でシステムの全体像を示した後，~\ref{sec:doa}-~\ref{sec:spec}で音響表現の生成手法を，~\ref{sec:network}でニューラルネットワークのアーキテクチャをそれぞれ詳述する．
	提案システムの概要をFig.~\ref{fig:fig1}に示す．

\subsection{問題設定}
\label{sec:problem}
	本研究では，ロボットマニピュレータが音を手がかりとして操作対象を識別し，把持して所定の位置へ搬送するタスクを扱う．各時刻 $t$ における観測 $o_t$ を入力とし，エキスパート行動 $a_t$ を再現するPolicy$\pi_\theta$を模倣学習によって獲得することを目的とする．データセットを $\mathcal{D}=\{(o_t,a_t)\}$ とすると，共通の目的は次の経験損失を最小化することである：
	$$
	\min_{\theta}\ \mathbb{E}_{(o_t,a_t)\sim\mathcal{D}}
	\left[\mathcal{L}_{\mathrm{BC}}\left(\pi_\theta(o_t),a_t\right)\right],
	$$
	ここで $\theta$ はPolicyの学習パラメータ，$\mathbb{E}_{(o_t,a_t)\sim\mathcal{D}}[\cdot]$ はデータセット $\mathcal{D}$ からサンプルした観測・行動ペアに関する期待値，$\pi_\theta(o_t)$ はPolicyが観測 $o_t$ から予測する行動，$\mathcal{L}_{\mathrm{BC}}$ はエキスパート行動とのずれを測る模倣学習損失であり，具体的な形はPolicyごとに異なる．
	ACTでは行動列のL1誤差に潜在変数のKLダイバージェンスを加え，Diffusion Policyでは拡散過程で加えたノイズの平均二乗誤差を用いる．
	VQ-BeTでは離散化された行動コードの分類損失と連続オフセットのL1損失を用い，$\pi_0$ではflow matchingにおける速度場の平均二乗誤差を用いる．
	各損失の詳細は~\ref{sec:network}で述べる．

	観測 $o_t$ は，正面・側面 2 台の RGB 画像，音響表現，および固有受容感覚から構成される．
	固有受容感覚とは，ロボット自身の状態を表す内部センサ情報を指す．本研究では，シミュレーションでは手先位置・姿勢とグリッパ開度，実機では右腕7関節角を用いる．
	
	シミュレーションでは固有受容感覚は手先位置・姿勢とグリッパ開度からなる9次元ベクトルであり，行動 $a_t$ はFranka Pandaの7関節位置目標と2つのフィンガ位置目標からなる9次元の位置制御指令である．

	特に本研究が対象とする難しさは，視覚情報だけでは操作対象を一意に特定できない状況である．すなわち，外観上同一の複数の物体が存在し，音によってのみ正しい操作対象を識別できる音環境マニピュレーション問題を中心的な課題として扱う．


\subsection{Acoustic-Scene Manipulation Tasks}
\label{sec:tasks}
	本研究では，音響情報が操作判断に果たす役割を体系的に検証するため，
	必要な音響的手がかりの性質が異なる4種類の操作タスクを設計し，
	これらを総称して\textbf{Acoustic-Scene Manipulation Tasks}と呼ぶ．
	各タスクの概要を Table~\ref{tab:tasks} に示す．

	\textbf{SoundPosition}（音源定位タスク）は，

	マニピュレータの操作範囲内に視覚的に同一の2つの物体と，1つの箱が配置してある．
	視覚的に同一の2つの物体のうち，音を発している物体を把持して箱に入れるタスクである．

	2つの物体の音色は同一であり，
	判断には音源の空間的位置のみが必要となる．
	このタスクは，音空間マップが操作判断に寄与するかを検証するために設計した．

	\textbf{SoundTone}（音色識別タスク）は，
	1つの物体が発する音の種類を識別し，音色に応じて2種類の箱のいずれかへ搬送するタスクである．
	音源の位置は固定されているため，判断に空間情報は不要であり，
	スペクトログラムが音色の弁別に有効かを検証する．

	\textbf{SoundAll}（統合判断タスク）は，
	視覚的に同一の2つの物体のうち，特定の音色を発する方を把持し，
	音色に対応する箱へ搬送するタスクである．
	対象物体の選択に空間情報が，搬送先の決定に周波数情報がそれぞれ必要であり，
	音空間マップとスペクトログラムの統合的な活用を検証する，
	本研究で最も総合的なタスクである．

	\textbf{SoundShake}（能動的探索タスク）は，
	視覚的に同一の2つの物体を順に把持して振り，
	音が鳴った方の物体を箱に入れるタスクである．
	ロボットが自ら探索行動を行って初めて音響情報を得られる点，
	および把持後の音響フィードバックに基づく時間的な条件分岐
	（音が鳴ればそのまま搬送し，鳴らなければ置き直して別の物体を把持する）
	が求められる点で，他の3タスクと質的に異なる．

	\begin{table*}[t]
		\centering
		\scriptsize
		\setlength{\tabcolsep}{3pt}
		\caption{Acoustic-Scene Manipulation Tasks の概要と各タスクが要求する音響的能力．}
		\label{tab:tasks}
		\begin{tabular}{p{0.10\textwidth} p{0.54\textwidth} >{\centering\arraybackslash}p{0.09\textwidth} >{\centering\arraybackslash}p{0.09\textwidth} >{\centering\arraybackslash}p{0.09\textwidth}}
			\hline
			タスク & 説明 & 音の位置情報 & 音の種類識別 & 探索的行動 \\
			\hline
			SoundPosition & 視覚的に同一の2物体のうち，音が鳴っている方を把持して箱に入れる & $\checkmark$ & -- & -- \\
			SoundTone & 1つの物体が発する音の種類を識別し，2種類の箱へ分別する & -- & $\checkmark$ & -- \\
			SoundAll & 視覚的に同一の2物体のうち，特定の音色を発する方を把持して対応する箱に入れる & $\checkmark$ & $\checkmark$ & -- \\
			SoundShake & 視覚的に同一の2物体を順に把持して振り，音が鳴る方を箱に入れる & $\checkmark$ & -- & $\checkmark$ \\
			\hline
		\end{tabular}
	\end{table*}

\subsection{システム概要}
\label{sec:overview}
	提案システムは，音響情報取得・表現モジュールとマルチモーダル模倣学習モデルから構成される（Fig.~\ref{fig:fig1}）．作業空間には $n$ 個のマイクロフォンアレイと2台のRGBカメラを配置する．各アレイは半径3.5cmの円周上に均等配置された8素子からなる．音響情報取得モジュールは，多チャンネル音響信号から空間的な音源尤度を表す音空間マップと，音源の周波数特性を表すスペクトログラムを生成し，これらを視覚特徴とともにPolicyへ入力する．RGB画像・音空間マップ・スペクトログラムはすべて $224 \times 224$ の解像度で統一する．

\subsection{音空間マップの生成}
\label{sec:doa}
\subsubsection{DOA推定}
	各マイクロフォンアレイの多チャンネル信号に対し，MUSIC（MUltiple SIgnal Classification）法~\citep{schmidt1986music}を適用する．入力音声信号はサンプリング周波数16kHzで取得し，FFT長512，50\%オーバーラップの短時間フーリエ変換（STFT）で周波数領域へ変換する．$c\in\{1,\ldots,n\}$ をアレイのインデックス，$M$ を各アレイのマイクロフォン数とする．$c$ 番目のアレイで観測された $M$ チャンネル信号のSTFTを $\mathbf{x}_{c,f,\tau}\in\mathbb{C}^{M}$ とすると，周波数ビン $f$ における空間共分散行列は
	$$
	\mathbf{R}_{c,f}=\frac{1}{T}\sum_{\tau=1}^{T}
	\mathbf{x}_{c,f,\tau}\mathbf{x}_{c,f,\tau}^{\mathrm{H}}
	$$
	で与えられる．ここで $\mathbb{C}^{M}$ は $M$ 次元複素ベクトル空間，$\mathbf{R}_{c,f}$ は $c$ 番目のアレイにおける周波数ビン $f$ の空間共分散行列，$T$ はSTFTの時間フレーム数，$\tau$ は時間フレームのインデックス，$\mathrm{H}$ は複素共役転置を表す．
	MUSIC法は，$\mathbf{R}_{c,f}$ の固有値分解から雑音部分空間 $\mathbf{E}_{n,c,f}$ を求め，方向 $\theta$ のステアリングベクトル $\mathbf{a}_{c,f}(\theta)$ との直交性に基づいて空間スペクトル
	$$
	P_c(\theta)=\frac{1}{|\mathcal{F}|}\sum_{f\in\mathcal{F}}
	\frac{1}{\mathbf{a}_{c,f}(\theta)^{\mathrm{H}}
	\mathbf{E}_{n,c,f}\mathbf{E}_{n,c,f}^{\mathrm{H}}
	\mathbf{a}_{c,f}(\theta)+\epsilon}
	$$
	を計算する．ここで $\mathbf{E}_{n,c,f}$ は雑音部分空間を張る固有ベクトルを列に持つ行列，添字 $n$ はnoise subspaceを表す．$\mathcal{F}$ は平均化に用いる周波数ビン集合，$|\mathcal{F}|$ はその要素数，$\epsilon$ はゼロ除算を避けるための微小正数である．
	シミュレーションでは推定音源数を $K_{\rm src}=3$ とした．方位ごとの空間スペクトルを対数化および正規化した値をDOAスコアとして用いる．

\subsubsection{二次元マップへの投影}
	推定されたDOAスコアを作業空間に対応した $224\times224$ の二次元平面へ投影し，$n$ チャンネルの音空間マップを生成する．画像の行インデックスを $i$，列インデックスを $j$ とし，画像座標 $u=(i,j)$ に対応する作業空間上の点を $p(u)$，アレイ中心を $m_c$ とすると，各チャンネルの値は
	$$
	S_c(u)=\frac{\bar{P}_c(\theta(p(u),m_c))}
	{\left(\|p(u)-m_c\|_2+d_{\mathrm{floor}}\right)^\gamma}
	$$
	で計算する．ここで $\bar{P}_c$ は正規化後のDOAスコア，$\theta(p(u),m_c)$ は $m_c$ から $p(u)$ へ向かう方位角，$\|\cdot\|_2$ はユークリッド距離，$d_{\mathrm{floor}}$ と $\gamma$ は距離減衰を安定化するパラメータである．
	シミュレーションでは $d_{\mathrm{floor}}=10.0\,\mathrm{m}$，$\gamma=1.0$ とした．
	全アレイの $S_c$ を重ねた多チャンネルマップ $S\in\mathbb{R}^{224\times224\times n}$ を空間音響表現として用いる．ここで $\mathbb{R}$ は実数の集合である．
	Policyへ入力する音空間マップは $S$ の $n$ チャンネルであり，本研究の標準設定では $n=4$ とする．
	この表現は，単一の推定座標に情報を圧縮せず，各アレイから見た方向尤度を保持するため，視覚特徴との空間的な対応づけを学習しやすい．

\subsection{スペクトログラムの抽出}
\label{sec:spec}
\subsubsection{Spotformingによる音源強調}
	音空間マップだけでは，音色の違いを捉えることが困難である．そこでSpotforming技術~\citep{spotforming_nmf}に基づき，目的音源方向の強調とNMF~\citep{lee1999nmf}による共通成分抽出を組み合わせ，雑音やサイドローブを抑制したスペクトログラムを生成する．まず全アレイのマップを
	$$
	S_{\mathrm{sum}}(u)=\sum_{c=1}^{n} S_c(u)^\alpha
	$$
	として統合し，ガウシアンフィルタ（$\sigma=1.0$）を適用した後，局所最大点を音源候補として選択する（本研究では最大1候補を使用）．ここで $S_{\mathrm{sum}}$ は全アレイを統合した音源候補マップ，$\alpha$ は各アレイのスコアを強調する指数であり，シミュレーションでは $\alpha=1.0$ とした．選択された音源候補を各アレイから見た方位を $\hat{\theta}_c$ とし，各アレイでは $\hat{\theta}_c$ に対してDelay-and-Sumビームフォーミングを行う：
	$$
	y_c(t)=\frac{1}{M}\sum_{m=1}^{M}
	x_{c,m}(t-\tau_{c,m}(\hat{\theta}_c)).
	$$
	ここで $t$ は時刻，$x_{c,m}(t)$ は $c$ 番目のアレイの $m$ 番目のマイクロフォン信号，$\tau_{c,m}(\hat{\theta}_c)$ は方位 $\hat{\theta}_c$ から到来する音が各マイクに届くまでの相対遅延である．
	次に各アレイのビームフォーミング出力 $y_c$ の振幅スペクトログラムを時間方向に結合した非負値行列 $\mathbf{V}\in\mathbb{R}_{\ge0}^{F_b\times nT_b}$ を作る．ここで $F_b$ は周波数ビン数，$T_b$ は各アレイの時間フレーム数である．$\mathbf{V}$ を50成分のNMFにより
	$$
	\mathbf{V}\approx \mathbf{W}\mathbf{H},\quad
	\mathbf{W}\ge0,\ \mathbf{H}\ge0
	$$
	と分解する．ここで $\mathbb{R}_{\ge0}$ は非負実数の集合，$\mathbf{W}\in\mathbb{R}_{\ge0}^{F_b\times R}$ は基底行列，$\mathbf{H}\in\mathbb{R}_{\ge0}^{R\times nT_b}$ は活性行列，$R=50$ はNMF成分数である．
	活性行列をアレイごとの時間区間に分割したものを $\mathbf{H}^{(c)}\in\mathbb{R}_{\ge0}^{R\times T_b}$ とし，各アレイで共通して強い成分を
	$$
	\mathbf{Z}=\min_{c}\mathbf{H}^{(c)}
	$$
	として推定する．ここで $\min_c$ は各NMF成分・各時刻についてアレイ方向の最小値を取る演算である．
	$\mathbf{Z}$ を最大値で正規化した行列を $\tilde{\mathbf{Z}}$ とし，閾値 $\delta=1.6\times10^{-3}$ によるバイナリマスク $\mathbf{B}$ を
	$$
	\mathbf{B}_{r,\tau}=
	\begin{cases}
	1 & \tilde{\mathbf{Z}}_{r,\tau}>\delta,\\
	0 & \mathrm{otherwise}
	\end{cases}
	$$
	として作成する．ここで $r$ はNMF成分のインデックス，$\tau$ はNMF活性行列における時間フレームのインデックス，$\mathbf{B}_{r,\tau}=1$ はその成分を再構成に残すことを表す．
	本研究ではマスクが広がりすぎることを避けるため，$\mathbf{B}$ の非ゼロ比率が0.2を超える場合は上位20\%の活性のみを残すように閾値を適応的に引き上げる．
	最後に $\mathbf{W}(\mathbf{H}^{(c)}\odot\mathbf{B})$ により各アレイの振幅スペクトログラムを再構成する．ここで $\odot$ は要素ごとの積である．

	位相には同じアレイのビームフォーミング出力 $y_c$ のSTFT位相を用い，逆STFTで得た波形をアレイ間で平均する．この処理は，目的音源が全アレイで観測される一方，サイドローブや局所雑音は一部のアレイでのみ強く現れるという仮定に基づく．

\subsubsection{対数パワースペクトログラムの生成}
	NMFで抽出した共通成分を再構成・統合し，雑音を抑制した単一チャンネルの対数パワースペクトログラムを生成する．再構成波形 $\hat{y}(t)$ のSTFTを $\hat{Y}_{f,\tau}$ とすると，スペクトログラム画像は
	$$
	G(f,\tau)=10\log_{10}\left(|\hat{Y}_{f,\tau}|^2+\epsilon\right)
	$$
	として計算する．ここで $f$ は周波数ビン，$\tau$ は時間フレーム，$\epsilon=10^{-10}$ は対数計算を安定化する微小正数である．
	表示周波数帯域と画像化時の正規化はシミュレーション評価と実機評価で異なり，具体的な値はAppendix~\ref{sec:appendixA4}に示す．
	正規化後のスペクトログラムは $224\times224$ にリサイズしてPolicyへ入力する．
	Policyに入力するスペクトログラムは1チャンネルである．
	得られた $G$ を周波数音響表現として用いることで，音空間マップでは失われる音色や周期性の情報をPolicyへ与える．

\subsection{ネットワークアーキテクチャ}
\label{sec:network}
	本研究では行動生成Policyとして\textbf{ACT}~\citep{act}，\textbf{Diffusion Policy}~\citep{diffusionpolicy}，\textbf{VQ-BeT}~\citep{vqbet}，\textbf{$\pi_0$}~\citep{pi0} の4種を評価する．
	本研究のアーキテクチャ上の変更点は，各Policyの行動生成部を大きく変更するのではなく，音空間マップとスペクトログラムを既存Policyが扱えるトークンまたは条件ベクトルへ変換する入力エンコーダを追加した点である．
	標準設定では，正面RGB画像，側面RGB画像，$n$ チャンネル音空間マップ，1チャンネルスペクトログラム，固有受容感覚を入力する．
	音空間マップは $n$ 個のマイクロフォンアレイに対応するチャンネルを持ち，本実験の標準設定では $n=4$ である．
	スペクトログラムはすべてのPolicyで1チャンネル入力として扱う．
	主要な学習パラメータを本文中の~\ref{sec:sim_setup}とAppendix~\ref{sec:appendixA}に示す．

	\begin{description}
		\item[\textbf{視覚エンコーダ}] 正面・側面2枚のRGB画像を視覚エンコーダへ入力し，物体位置，箱位置，ロボット姿勢に関する特徴量を抽出する．
		\item[\textbf{音空間マップエンコーダ}] $n$ チャンネルの音空間マップを専用の畳み込みエンコーダへ入力し，音源の空間的な尤度分布を特徴化する．RGB画像とは入力チャンネル数と統計が異なるため，音空間マップ用の最初の畳み込み層は $n$ チャンネル入力として初期化し，ImageNet事前学習重みは用いない．
		\item[\textbf{スペクトログラムエンコーダ}] 対数パワースペクトログラムを1チャンネル入力の専用畳み込みエンコーダへ入力し，音色や周期性に対応する周波数パターンを抽出する．このエンコーダもRGB画像とは別の重みを持つ．
	\end{description}

\subsubsection{ACT}
	ACTは，複数時刻分の行動を一度に出力するAction Chunking Transformerである．
	Action Chunkとは，時刻 $t$ で1ステップ先の行動だけでなく，$a_t,\ldots,a_{t+H-1}$ のような長さ $H$ の未来行動列をまとめて予測する単位を指す．本実験では $H=100$ とした．
	正面・側面RGB画像にはImageNet事前学習済みResNet-18~\citep{he2016resnet}を用い，音空間マップとスペクトログラムには別々のResNet-18を追加する．音空間マップ用ResNetの入力チャンネル数は $n=4$，スペクトログラム用ResNetの入力チャンネル数は1である．
	各ResNetの最終特徴マップを $1\times1$ 畳み込みでTransformerの隠れ次元512へ射影し，2次元位置埋め込みを加えたうえでトークン列へ変換する．
	音空間マップは空間的な位置対応を保つため，Global Average Poolingで1ベクトルに潰さず，特徴マップ上の各位置をトークンとしてTransformer encoder~\citep{vaswani2017attention}へ入力する．

	ACTの訓練時には，行動列から潜在変数 $z$ を推定するVAEを併用する．VAEとはVariational Autoencoderの略で，訓練時に行動列を低次元の確率変数へ圧縮し，その分布を標準正規分布へ近づけることで多様な行動列を扱いやすくする仕組みである．
	訓練損失は
	$$
	\mathcal{L}_{\mathrm{ACT}}
	=\frac{1}{H d_a}\sum_{h=0}^{H-1}
	m_h\left\|\hat{a}_{t+h}-a_{t+h}\right\|_1
	+\beta_{\mathrm{KL}}\,
	D_{\mathrm{KL}}\left(q_\phi(z\mid a_{t:t+H-1},s_t)\,\|\,\mathcal{N}(0,\mathbf{I})\right)
	$$
	である．ここで $d_a$ は行動次元，$h$ は行動チャンク内の時刻インデックス，$\hat{a}_{t+h}$ は予測行動，$a_{t+h}$ はエキスパート行動，$m_h\in\{0,1\}$ はデータ列端のパディングを損失から除外するためのマスク，$s_t$ は固有受容感覚，$z$ はVAEの潜在変数，$\phi$ はVAE encoderのパラメータ，$q_\phi$ はVAE encoderが出力する潜在分布，$\mathcal{N}(0,\mathbf{I})$ は平均0・共分散単位行列の標準正規分布，$D_{\mathrm{KL}}$ はKLダイバージェンス，$\beta_{\mathrm{KL}}=10$ はKL項の重みである．
	評価時のシミュレーションでは，ACTのみ時間アンサンブルを用い，1ステップごとにチャンクを再予測して指数重み係数0.01で同一時刻に対応する複数の予測行動を平均した．

\subsubsection{Diffusion Policy}
	Diffusion Policyでは，行動列をノイズから逐次的に復元する条件付き拡散モデル~\citep{ho2020ddpm, diffusionpolicy}を用いる．
	拡散モデルとは，訓練時には正解行動列へ段階的にガウスノイズを加え，推論時にはその逆過程を学習済みネットワークでたどる生成モデルである．
	RGB画像，音空間マップ，スペクトログラムをそれぞれResNet-18~\citep{he2016resnet}とSpatial Softmaxで1次元特徴へ変換する．
	Spatial Softmaxとは，特徴マップ上の各位置を確率分布として扱い，特徴が強く反応した位置の期待座標を取り出す操作であり，本実験では32個のキーポイントを用いる．
	音空間マップとスペクトログラムにはACTと同様に専用ResNetを追加し，RGB画像用ResNetとは重みを共有しない．
	2ステップ分の観測履歴から得た画像特徴と固有受容感覚を結合し，U-Netのグローバル条件として入力する．グローバル条件とは，U-Netの全時刻に共通して与える条件ベクトルであり，ここでは現在と直前の観測情報をまとめたものである．

	訓練では，長さ $H_d=16$ の行動列 $x_0=a_{t-1:t+14}$ に対し，100段階の拡散ステップから一様にサンプルした $k$ でノイズ $\epsilon\sim\mathcal{N}(0,\mathbf{I})$ を加えた
	$$
	x_k=\sqrt{\bar{\alpha}_k}x_0+\sqrt{1-\bar{\alpha}_k}\epsilon
	$$
	を作る．ここで $\bar{\alpha}_k$ はDDPMのノイズスケジュールから定まる係数である．U-Netを $\epsilon_\theta$，観測条件を $g(o)$ とすると，損失は
	$$
	\mathcal{L}_{\mathrm{Diff}}
	=\mathbb{E}_{x_0,\epsilon,k}
	\left[\frac{1}{H_d d_a}\left\|\epsilon_\theta(x_k,k,g(o))-\epsilon\right\|_2^2\right]
	$$
	である．ここで $\theta$ はU-Netの学習パラメータ，$o$ は観測，$g(o)$ は観測をエンコードした条件ベクトル，$d_a$ は行動次元である．推論時は100ステップの逆拡散を行い，観測履歴の最後の時刻から8ステップ分の行動を実行候補として取り出す．

\subsubsection{VQ-BeT}
	VQ-BeTは，連続的な行動列を離散コードと連続オフセットの組として予測するBehavior Transformerである．
	離散コードとは，似た行動チャンクを有限個の代表パターンへ割り当てるインデックスであり，連続オフセットとは，その代表パターンから正解行動へ微調整するための残差ベクトルである．
	RGB画像，音空間マップ，スペクトログラムに専用ResNet-18~\citep{he2016resnet}を用意し，各特徴マップをSpatial Softmaxで1次元特徴へ変換する．
	音空間マップ用ResNetは $n$ チャンネル入力を処理し，スペクトログラム用ResNetは1チャンネル入力を処理する．
	得られた各モダリティ特徴，固有受容感覚，およびaction query tokenをGPTスタイルのpolicy networkへ入力する．
	action query tokenとは，画像や状態を表す観測トークンではなく，「この位置に対応する行動チャンクを出力せよ」という問い合わせ用の学習可能トークンである．
	本実験では観測履歴長を5，1つのaction query tokenが予測する行動チャンク長を5，現在チャンクに加えて2個の未来チャンク用query tokenを追加した．

	訓練は2段階で行う．第1段階ではResidual VQ-VAEを20,000ステップ学習し，長さ5の行動チャンクを2層のResidual Vector Quantization（RVQ）コードへ離散化する．
	RVQとは，1つ目のコードで粗い代表ベクトルを選び，2つ目以降のコードで残差を順に量子化する方法である．
	codebookとは離散コードが参照する代表ベクトルの集合であり，本研究では各RVQ層のcodebook sizeを16とした．
	第2段階では，固定したVQ-VAEに対して，policy networkがコード確率とオフセットを予測する．
	予測行動チャンクを $\hat{A}_j$，正解行動チャンクを $A_j$，1層目と2層目の正解コードをそれぞれ $c_j^{(1)},c_j^{(2)}$，予測コード分布を $p_{\theta,j}^{(1)},p_{\theta,j}^{(2)}$ とすると，損失は
	$$
	\mathcal{L}_{\mathrm{VQ}}
	=\frac{1}{J}\sum_{j=1}^{J}\left[
	\lambda_{\mathrm{off}}\frac{1}{H_q d_a}\|\hat{A}_j-A_j\|_1
	+\lambda_1\,\mathrm{FL}(p_{\theta,j}^{(1)},c_j^{(1)})
	+\lambda_2\,\mathrm{FL}(p_{\theta,j}^{(2)},c_j^{(2)})
	\right]
	$$
	である．ここで $\theta$ はpolicy networkの学習パラメータ，$j$ はaction query tokenのインデックス，$J=3$ はaction query token数，$H_q=5$ は1つのaction query tokenが担当する行動チャンク長，$d_a$ は行動次元である．
	$\mathrm{FL}$ はFocal Lossであり，正解コード $c$ に割り当てられた確率を $p_c$ とすると
	$$
	\mathrm{FL}(p,c)=-(1-p_c)^2\log p_c
	$$
	で定義される．この損失は，分類が容易なサンプルの寄与を小さくし，難しいコード分類を重視する．本研究では $\lambda_{\mathrm{off}}=10000$，$\lambda_1=5.0$，$\lambda_2=0.5$ とした．

\subsubsection{$\pi_0$}
	$\pi_0$は，PaliGemma~\citep{beyer2024paligemma}の視覚言語モデルと行動生成用のGemma expertを組み合わせ，flow matching~\citep{lipman2023flow}で行動列を生成するPolicyである．
	prefixとは，行動生成の前提条件としてモデルに先に与えられるトークン列を指し，本研究では画像トークン，音響トークン，言語トークンをprefixとして扱う．
	suffixとは，prefixを条件として生成・更新される側のトークン列を指し，ここでは固有受容感覚，ノイズ付き行動，時刻埋め込みから構成される．
	正面・側面RGB画像はPaliGemmaの画像エンコーダでトークン化する．一方，音空間マップとスペクトログラムはPaliGemmaの画像エンコーダへ直接入力せず，専用ResNet-18で特徴マップ化する．その後，$1\times1$ 畳み込みでPaliGemmaの隠れ次元へ射影し，2次元位置埋め込みを加えてprefixへ追加する．
	この設計により，RGB画像に対する事前学習済み表現を保ちつつ，音空間マップとスペクトログラムを異なる統計を持つ画像状モダリティとして扱える．

	flow matchingとは，ノイズ分布からデータ分布へ移動する連続的な速度場を学習する生成手法である．正解行動列を $a$，同じ形状の標準正規ノイズを $\epsilon$，時刻を $r=0.999b+0.001$，$b\sim\mathrm{Beta}(1.5,1.0)$ とすると，訓練時には
	$$
	x_r=r\epsilon+(1-r)a,\qquad u_r=\epsilon-a
	$$
	を作り，モデル $v_\theta$ が速度 $u_r$ を予測するように
	$$
	\mathcal{L}_{\pi_0}
	=\mathbb{E}_{a,\epsilon,r}
	\left[\frac{1}{H_\pi d_a}\left\|v_\theta(x_r,r,o)-u_r\right\|_2^2\right]
	$$
	を最小化する．ここで $\theta$ は速度場モデルの学習パラメータ，$H_\pi=50$ は行動チャンク長，$d_a$ は行動次元，$x_r$ はノイズ付き行動，$u_r$ はノイズから正解行動へ向かう教師速度，$o$ はprefixと固有受容感覚を含む観測条件である．
	本実験ではPaliGemma側に事前学習済み重みを用い，音響エンコーダは追加モジュールとして学習した．

%===============================================================================
\section{Experiments}
\label{sec:result}
	本節では，提案する音響表現が模倣学習Policyの性能に与える影響を，以下の3点において検証する．
	(1)~音空間マップとスペクトログラムが有効なタスク
	(2)~タスクに不要な音響モダリティを入力に加えた場合の性能変化
	(3)~各Policyにおける音響情報の活用効率の差異
	まず~\ref{sec:sim_setup}でシミュレーション環境と評価手順を述べ，~\ref{sec:sim_results}で定量・定性的な実験結果を報告し，~\ref{sec:real}で実機実験による検証を行う．

\subsection{実験設定}
\label{sec:sim_setup}
	~\ref{sec:tasks}で定義した4タスクについて，各音響表現の寄与を切り分けて評価するため，
	全Policyに対して以下の4条件を設定し，評価を行った．

	\begin{description}
		\item[\textbf{Baseline}] 視覚情報のみ（音響情報なし）．
		\item[\textbf{Ours$_{V,Spa}$}] 視覚情報 + 音空間マップ．
		\item[\textbf{Ours$_{V,Spec}$}] 視覚情報 + スペクトログラム．
		\item[\textbf{Ours$_{V,Spa,Spec}$}] 視覚情報 + 音空間マップ + スペクトログラム（提案手法）．
	\end{description}

	物理シミュレータGenesis~\citep{Genesis}と音響シミュレータPyroomacoustics~\citep{pyroomacoustics}を統合したシミュレーション環境を構築した（Fig.~\ref{fig:fig2}）．ロボットアームにはFranka Emika Pandaを使用し，行動次元は9とした．視覚情報は正面・側面カメラから取得した $224 \times 224$ のRGB画像2枚とした．音響シミュレーションでは $10\,\mathrm{m}\times10\,\mathrm{m}\times3\,\mathrm{m}$ の室内環境を用い，最大3次反射まで考慮した．マイクロフォンアレイ数は予備実験（Appendix~\ref{sec:appendixA1}）に基づき，標準設定で $n=4$ とし，作業空間中心から半径300mmの円周上に配置した．視覚情報は30FPSで更新され，音情報は予備実験（Appendix~\ref{sec:appendixA2}）に基づき，3フレームごと（10FPS）で更新した．

	上記のシミュレーション環境において，デモンストレーションデータを収集した．
	デモンストレーションはイベントベースでエンドエフェクタの目標Poseを指定し，IKベースのエキスパートにより軌道を生成することにより収集した．
	
	SoundPosition，SoundToneタスクについては100エピソード，SoundAll，SoundShakeタスクについては200エピソードのデモンストレーションを収集した．
	学習設定として，SoundPositionおよびSoundToneタスクでは10万ステップ，SoundAllおよびSoundShakeタスクでは 20 万ステップの学習を行った．
	バッチサイズは$\pi_0$で4，ACTで8，Diffusion PolicyとVQ-BeTで32とした．
	学習率はACTで $1.0\times10^{-5}$，Diffusion PolicyとVQ-BeTで $1.0\times10^{-4}$，$\pi_0$で $2.5\times10^{-5}$ とした．
	VQ-BeTのResidual VQ-VAE事前学習段階のみ，学習率 $1.0\times10^{-3}$ で20,000ステップ学習した．
	学習済み重みは，10万ステップ学習の条件では10,000ステップ間隔，20万ステップ学習の条件では20,000ステップ間隔で評価した．
	これらのバッチサイズ，学習率，学習ステップ数，行動チャンク長は性能に大きく影響するため本文中に明記し，その他の最適化設定とモデル内部パラメータはAppendix~\ref{sec:appendixA}にまとめる．
	各タスクは3つの異なる乱数シードで学習し，評価指標は100回の試行において700ステップ以内にタスクを完遂した割合とした．
	なお，各タスクの成功率は把持・搬送・投入を含む全体完遂率であり，対象物体の選択正答率とは区別される．
	学習曲線をFig.~\ref{fig:fig3}に示す．

	\begin{figure}[t]
		\centering
		\includegraphics[keepaspectratio, width=14cm]{fig/fig2.png}
		\caption{Simulation settings. (A) 物理シミュレータGenesisと音響シミュレータPyroomacousticsを統合したシステム構成．Franka Emika Pandaの周囲に4つのマイクロフォンアレイを配置．(B) 4種類の評価タスク（SoundPosition，SoundTone，SoundAll，SoundShake）．}
		\label{fig:fig2}
	\end{figure}

\subsection{シミュレーション実験結果}
	全タスク・全Policy・全条件の成功率をTable~\ref{tab:results}に示す．また，学習ステップ数に対する成功率の推移をFig.~\ref{fig:fig3}に示す．

	\begin{figure}[t]
		\centering
		\includegraphics[keepaspectratio, width=14cm]{fig/fig3.png}
		\caption{学習ステップ数に対するタスク成功率の変化．4タスク（SoundPosition，SoundTone，SoundAll，SoundShake）における4 Policy（ACT，Diffusion Policy，VQ-BeT，$\pi_0$）のOurs3条件での学習推移を示す．}
		\label{fig:fig3}
	\end{figure}

	\begin{table*}[t]
		\centering
		\scriptsize
		\setlength{\tabcolsep}{4pt}
		\caption{全タスクにおけるタスク成功率（\%）の比較（Mean $\pm$ STD）．各Policy内で最良の値を太字で示す．}
		\label{tab:results}
		\begin{tabular}{llcccc}
			\hline
			モデル & 条件 & SoundPosition & SoundTone & SoundAll & SoundShake \\
			\hline
			ACT & Baseline & 20.3 $\pm$ 7.2 & 39.3 $\pm$ 4.0 & 11.0 $\pm$ 3.5 & 35.0 $\pm$ 5.3 \\
			& Ours1 & \textbf{74.7 $\pm$ 4.7} & 44.3 $\pm$ 2.3 & 45.0 $\pm$ 3.6 & 75.3 $\pm$ 5.0 \\
			& Ours2 & 15.3 $\pm$ 4.0 & 74.3 $\pm$ 7.0 & 25.7 $\pm$ 3.1 & 72.3 $\pm$ 8.4 \\
			& Ours3 & 73.0 $\pm$ 5.6 & \textbf{76.7 $\pm$ 1.2} & \textbf{78.7 $\pm$ 2.1} & \textbf{77.3 $\pm$ 5.7} \\
			\hline
			Diffusion Policy & Baseline & 13.3 $\pm$ 1.2 & 41.3 $\pm$ 2.5 & 9.3 $\pm$ 6.1 & 26.3 $\pm$ 0.6 \\
			& Ours1 & \textbf{68.3 $\pm$ 13.7} & 32.7 $\pm$ 3.1 & 46.0 $\pm$ 5.3 & 37.3 $\pm$ 8.1 \\
			& Ours2 & 15.3 $\pm$ 1.2 & \textbf{80.3 $\pm$ 3.8} & 8.7 $\pm$ 3.8 & 6.3 $\pm$ 2.3 \\
			& Ours3 & 68.0 $\pm$ 3.6 & 67.7 $\pm$ 7.2 & \textbf{89.7 $\pm$ 2.1} & \textbf{44.0 $\pm$ 3.6} \\
			\hline
			VQ-BeT & Baseline & 12.7 $\pm$ 6.8 & 24.0 $\pm$ 6.9 & 7.7 $\pm$ 1.2 & 17.3 $\pm$ 5.8 \\
			& Ours1 & \textbf{44.7 $\pm$ 5.5} & 32.7 $\pm$ 3.1 & 32.3 $\pm$ 7.5 & 20.0 $\pm$ 20.0 \\
			& Ours2 & 12.3 $\pm$ 3.8 & \textbf{49.3 $\pm$ 9.8} & 19.0 $\pm$ 1.0 & 25.0 $\pm$ 3.6 \\
			& Ours3 & 40.3 $\pm$ 4.5 & 32.3 $\pm$ 10.0 & \textbf{55.3 $\pm$ 11.7} & \textbf{32.7 $\pm$ 12.7} \\
			\hline
			$\pi_0$ & Baseline & 18.3 $\pm$ 2.3 & 32.3 $\pm$ 2.1 & 7.0 $\pm$ 4.0 & 9.0 $\pm$ 1.0 \\
			& Ours1 & 17.3 $\pm$ 5.0 & 40.0 $\pm$ 11.4 & 14.3 $\pm$ 2.9 & \textbf{26.7 $\pm$ 3.2} \\
			& Ours2 & 14.7 $\pm$ 5.7 & \textbf{78.7 $\pm$ 6.0} & 19.7 $\pm$ 3.8 & 9.3 $\pm$ 10.4 \\
			& Ours3 & \textbf{32.0 $\pm$ 9.9} & 73.3 $\pm$ 8.4 & \textbf{45.7 $\pm$ 8.1} & 18.7 $\pm$ 4.2 \\
			\hline
		\end{tabular}
	\end{table*}

\subsubsection{SoundPositionタスク：音源定位}
	SoundPositionタスクでは，音源の空間位置情報のみが判断に必要であり，音色情報は不要である．この特性はTable~\ref{tab:results}の結果に明確に反映されている．ACTおよびDiffusion Policyにおいて，Ours1はそれぞれ74.7\%・68.3\%と高い成功率を達成し，Ours2はBaselineとほぼ同水準にとどまった．Ours3も73.0\%・68.0\%を達成し，音空間マップが音源定位に有効な入力表現であることを示した．

	Baselineの成功率はいずれのPolicyでも20\%以下であった．本タスクの対象選択は二択であるが，Table~\ref{tab:results}の値は選択正答率ではなく，把持・搬送・箱投入までを含む全体成功率であるため，単純なランダム選択の期待値50\%とは直接比較できない．それでも，視覚的に対象を区別できない条件では全体成功率が大きく低下しており，誤選択に加えて，曖昧な入力に対して平均化された行動が生成され，把持段階で失敗したことを示している．

	$\pi_0$については，Ours3でも成功率が32.0\%にとどまり，他のPolicyと傾向が異なる．この結果は，大規模VLAモデルに音響特徴を統合する際には，特徴の挿入位置や学習設定が性能を大きく左右することを示している．

	なお，SoundPosition タスクでは Ours2（スペクトログラムのみ）が Baseline と同水準もしくはそれを下回るPolicyが見られた．本タスクの判断には音源の空間情報のみが必要であり，周波数情報は判断に寄与しない．デモンストレーション数が限られた条件下では，判断に不要なモダリティを入力に加えることが，Policyがそれを無視する表現を学習する妨げとなり，結果として性能を低下させる可能性がある．同様の傾向は SoundShake タスクの Diffusion Policy においても観察された（Ours2: 6.3\%，Baseline: 26.3\%）．これは，入力モダリティの選択がタスクの性質に応じて慎重に行われるべきことを示唆している．

\subsubsection{SoundToneタスク：音色識別}
	SoundTone タスクでは，音源位置は固定されており，音色の識別が主要な判断手がかりとなる．このタスクは視覚情報とスペクトログラムのみで原理的に解けるよう設計されている．π0 では Ours2（スペクトログラムのみ）が 78.7\% という最高水準の成功率を達成しており，スペクトログラム入力単独でも音色識別が学習可能であることを示している．一方，ACT・Diffusion Policy では Ours2 がそれぞれ 74.3\%，80.3\% と Baseline を大きく上回ったが，VQ-BeT では Ours2 が 49.3\% と他のPolicyに比べて低水準にとどまった．VQ-BeT は行動を離散コードに量子化する設計であり，音色に応じた連続的な搬送先調整との相性が悪い可能性がある．

	特に興味深いのは$\pi_0$の結果である．$\pi_0$はOurs2（スペクトログラムのみ）で78.7\%という最高水準の成功率を達成した．$\pi_0$はPaliGemmaをバックボーンとする事前学習済みの大規模モデルであり，スペクトログラムのような二次元マップ形式の入力表現と親和性が高い可能性がある．一方，音空間マップを追加したOurs3では成功率が73.3\%とやや低下しており，$\pi_0$においては固定された音源位置に関する空間情報が一部の事例でノイズとして作用した可能性が示唆される．

	Diffusion PolicyのOurs3は67.7\%と高い成功率を達成した一方，VQ-BeTは全条件で30\%台にとどまった．VQ-BeTは行動の離散化表現を用いるため，音色に応じた細かい搬送先の切り替えに必要な連続的な行動調整の表現が困難であった可能性がある．

\subsubsection{SoundAllタスク：音の位置と種類の統合}
	SoundAll タスクは，音源位置と音色の両方が必要な，本研究で最も総合的なタスクである．Baseline の成功率は全Policyで 7〜11\% と低く，視覚情報のみでは操作対象を特定できないことが確認される．Ours1（音空間マップのみ）は ACT で 45.0\%，Diffusion Policy で 46.0\% まで向上するが，音色に応じた箱選択には不十分である．Ours2（スペクトログラムのみ）も対象物体の選択が不安定となり，性能は限定的である．Ours3 は，ACT で 78.7\%，Diffusion Policy で 89.7\%，VQ-BeT で 55.3\%，π0 で 45.7\% を達成し，全Policyで他の条件を上回った．この結果は，位置と音色の両方が判断に必要となるタスクでは，対応する両方の入力表現をPolicyに与えることが必要であることを示している．

	Diffusion PolicyのOurs3が89.7\%と全実験中最高の成功率を達成したことは注目すべき結果である．Diffusion Policyが生成するマルチモーダルな行動分布は，音源の位置と音色を組み合わせた複合的な判断に伴う複数の行動候補を扱ううえで有効に働いた．

\subsubsection{SoundShakeタスク：探索的行動と音響フィードバック}
	SoundShakeタスクは，ロボット自身が物体を把持して振るという探索的行動によって音響情報を引き出し，その結果に基づいて行動を切り替える能動的知覚タスクである．把持した物体が音を発せばそのまま箱へ搬送し，音が鳴らなければ置き直して別の物体を選択するという時間的な条件分岐が求められる．ACTのOurs1とOurs2はそれぞれ75.3\%・72.3\%を達成し，Ours3は77.3\%と同等以上の成功率を示した．本タスクでは，振動後に音が観測された物体を選択するイベント的な手がかりが強く，ACTではいずれの音響表現もBaselineを大きく上回った．Baselineも35.0\%を示しているのは，視覚的・動力学的な把持動作自体はデモンストレーションから学習でき，音情報がなくても部分的に再現されるためである．

	他のPolicyでは，Diffusion PolicyのOurs3が44.0\%を達成した一方，VQ-BeT・$\pi_0$では成功率が低水準にとどまった．SoundShakeタスクが要求する時間的に複雑な条件分岐の処理において，音響特徴の統合方法と時系列表現の設計が性能を大きく左右したと考えられる．

\subsubsection{横断的分析}
\label{sec:cross_analysis}
	4タスクにわたる結果（Table~\ref{tab:results}）から，以下の3つの横断的知見を抽出する．

	\textbf{音響モダリティとタスク要求の対応．}
	各タスクにおいて，タスクが要求する音響的能力に対応したモダリティを入力に含む条件が最良の性能を示す傾向が一貫して観察された．
	SoundPositionでは音空間マップを含む Ours1/Ours3 が，SoundTone ではスペクトログラムを含む Ours2/Ours3 が，SoundAll では両者を含む Ours3 がそれぞれ最良であった．
	この結果は，音響表現の設計がタスク要求と整合する場合にPolicyの判断能力が向上することを示している．

	\textbf{不要モダリティの追加による性能低下．}
	タスクに不要なモダリティを追加すると，性能が Baseline 以下に低下する場合がある．
	SoundPosition における Ours2（ACT: 15.3\%，Baseline: 20.3\%）や，SoundShake における Diffusion Policy の Ours2（6.3\%，Baseline: 26.3\%）がその典型例である．
	限られたデモンストレーションデータのもとでは，判断に不要な入力次元を無視する表現を学習することが困難となり，	結果としてPolicyの最適化を阻害することが示唆される．

	\textbf{Policy間の音響情報活用効率の差異．}
	音響情報の追加による改善幅はPolicyによって大きく異なる．
	ACT と Diffusion Policy は多くのタスクで大きな改善を示した一方，VQ-BeT と $\pi_0$ では改善が限定的なタスクが複数存在した．
	VQ-BeT は行動の離散化に伴う表現力の制約が，$\pi_0$ は大規模事前学習モデルへの音響特徴の統合方法が，それぞれ制約要因となった可能性がある．
	この知見は，音響モダリティの統合が万能ではなく，Policyのアーキテクチャに応じた設計が必要であることを示している．

\subsubsection{潜在空間の可視化}
\label{sec:tsne}
	各Policyの内部表現が音響情報をどのように符号化しているかを定性的に分析するため，行動生成直前に得られる潜在特徴量を t-SNE~\citep{vandermaaten2008tsne}により2次元へ射影して可視化した．各モデルからの潜在状態の取り出し方および全Policyの結果はAppendix~\ref{sec:appendixB}に示す．ここでは，本研究の主張に最も関連する2つの観察をFig.~\ref{fig:fig4}にまとめる．

	\begin{figure}[t]
		\centering
		\includegraphics[keepaspectratio, width=14cm]{fig/fig4.png}
		\caption{潜在空間のt-SNE可視化．(a) Diffusion PolicyのSoundAllタスクにおける，Baseline，Ours1，Ours2，Ours3の比較．音源位置（Sound Coordinate）および音種（Sound Type）でラベル付けした．Ours3で最も明瞭な構造化が観察される．(b) SoundShakeタスクにおけるACTとDiffusion Policyの音源位置に関する潜在表現の比較．ACTは音源位置に応じた分離が明瞭である一方，Diffusion Policyでは分離が弱い．}
		\label{fig:fig4}
	\end{figure}

	\textbf{(a) Ours3における潜在空間の構造化（Diffusion Policy，SoundAll）．}
	Fig.~\ref{fig:fig4}(a)はDiffusion PolicyのSoundAllタスクにおける各条件の潜在空間を示す．Baselineではいずれのラベル（音源位置・音種）でも一貫した分離が見られない．Ours1では音空間マップが入力されるため，音源位置に対応した分離が観察される一方，音種の分離は不明瞭である．逆にOurs2ではスペクトログラムにより音種の分離が現れるが，音源位置の分離は限定的である．Ours3では音空間マップとスペクトログラムを併用した結果，音源位置と音種の双方に対応した構造化が確認できる．この観察は，SoundAllタスクにおいてOurs3が89.7\%という最高成功率を達成したという定量的結果（Table~\ref{tab:results}）と整合する．

	\textbf{(b) Policy間の差：SoundShakeにおける空間情報の保持．}
	Fig.~\ref{fig:fig4}(b)はSoundShakeタスクにおけるACTとDiffusion Policyの潜在空間を音源位置でラベル付けして示したものである．ACTでは音源位置に応じた明瞭な分離が見られる一方，Diffusion Policyでは分離が著しく弱い．SoundShakeタスクではACTのOurs3が77.3\%，Diffusion PolicyのOurs3が44.0\%とPolicy間で大きな性能差が生じているが，この差は，ACTが音源の空間情報を潜在表現として保持できているのに対し，Diffusion Policyでは音響的手がかりが潜在空間に十分に符号化されていないことに対応していると解釈できる．

	これらの観察は，Table~\ref{tab:results}の定量的結果と整合しており，各音響モダリティが担う役割の違い，およびPolicyごとの音響情報活用能力の違いを定性的に支持している．

\subsection{実機実験}
\subsubsection{実験設定}
	シミュレーションで得られた知見が実環境においても再現されるかを検証するため，ALOHA~\citep{act}タイプの双腕型オープンソースロボットILOHAと4台の円形マイクロフォンアレイTAMAGO-03を用いた実機実験環境を構築した．ILOHAの行動次元は14である．マイクロフォンアレイは，俯瞰視点で短辺60 cm，長辺120 cmの長方形の各頂点に配置した．
	実機実験は，シミュレーションにおいてOurs3とBaselineの差が最も顕著であったSoundAllタスクと，能動的探索を要するSoundShakeタスクの2種類について実施した．評価Policyには，シミュレーションにおいて両タスクで安定した性能を示したACTを採用し，Ours3とBaselineを比較した．各条件についてデモンストレーションデータは100エピソード収集し，学習はて学習に用いた100回の試行を実施した．

\subsubsection{実験結果}
	各条件のタスク成功率をTable~\ref{tab:real}に示す．SoundAll タスクにおいて，ACT の Ours3 は 72/100 の成功率を示した一方，Baseline は 5/100 にとどまった．SoundShake タスクにおいても，Ours3 は 50/100 を達成し，Baseline の 12/100 を上回った．これらの結果は，視覚情報のみでは解けないタスクに対して，提案手法がシミュレーションで観察された性能向上を実環境でも再現することを示している．
	実環境では，反射音，背景雑音，ロボット駆動系から生じる機械的振動など，シミュレーションでは完全に再現できない音響的外乱が存在する．それにもかかわらず，提案手法が実機上で機能した要因として，音空間マップが単一の推定座標ではなく方向尤度の分布として空間情報を保持していること，および Spotforming/NMF による前処理が複数アレイ間で共通して観測される成分のみを残すことで雑音と局所的なサイドローブを抑制していることが挙げられる．


	\begin{figure}[t]
		\centering
		\includegraphics[keepaspectratio, width=14cm]{fig/fig5.png}
		\caption{実機実験の実行例．上段：SoundAllタスクにおける把持・搬送の連続フレーム．下段：音空間マップの推移．}
		\label{fig:fig5}
	\end{figure}


\appendix
%===============================================================================
\section{パラメータ設定}
\label{sec:appendixA}
	各ポリシーの学習・モデルパラメータを Table~\ref{tab:hyperparams} に示す．
	本文では性能に直接影響する学習率，バッチサイズ，学習ステップ数を述べたため，本節では最適化の詳細，行動系列長，モダリティエンコーダの設定をまとめる．

	\begin{table*}[t]
		\centering
		\scriptsize
		\caption{各ポリシーの学習ハイパーパラメータ．}
		\label{tab:hyperparams}
		\begin{tabular}{p{0.19\textwidth} p{0.18\textwidth} p{0.20\textwidth} p{0.20\textwidth} p{0.17\textwidth}}
			\hline
			& ACT & Diffusion Policy & VQ-BeT & $\pi_0$ \\
			\hline
			オプティマイザ & AdamW & Adam & Adam & AdamW \\
			学習率 & $1.0\times10^{-5}$ & $1.0\times10^{-4}$ & $1.0\times10^{-4}$（VQ-VAE: $1.0\times10^{-3}$） & $2.5\times10^{-5}$ \\
			Weight decay & $1.0\times10^{-4}$ & $1.0\times10^{-6}$ & $1.0\times10^{-6}$（VQ-VAE: $1.0\times10^{-4}$） & $1.0\times10^{-2}$ \\
			学習率スケジューラ & なし & cosine，warmup 500 steps & VQ-BeT scheduler，warmup 500 steps，VQ-VAE 20,000 steps & cosine decay，warmup 1,000 steps，decay 30,000 steps，final LR $2.5\times10^{-6}$ \\
			バッチサイズ & 8 & 32 & 32 & 4 \\
			行動系列長 & 行動チャンク長100，評価時ACT時間アンサンブルでは1ステップずつ実行 & 行動系列長16，実行8ステップ & action query tokenあたり行動チャンク長5，action query token 3個 & 行動チャンク長50，実行50ステップ \\
			観測履歴長 & 1 & 2 & 5 & 1 \\
			視覚バックボーン & ResNet-18，ImageNet事前学習 & ResNet-18，事前学習なし，カメラごとに独立 & ResNet-18，事前学習なし & PaliGemma画像エンコーダ \\
			音空間マップエンコーダ & ResNet-18，$n$チャンネル入力，事前学習なし & ResNet-18，$n$チャンネル入力，事前学習なし & ResNet-18，$n$チャンネル入力，事前学習なし & ResNet-18，$n$チャンネル入力，$1\times1$ 畳み込みでPaliGemma次元へ射影 \\
			スペクトログラムエンコーダ & ResNet-18，1チャンネル入力，事前学習なし & ResNet-18，1チャンネル入力，事前学習なし & ResNet-18，1チャンネル入力，事前学習なし & ResNet-18，1チャンネル入力，$1\times1$ 畳み込みでPaliGemma次元へ射影 \\
			モデル次元 & Transformer次元512，head数8，encoder 4層，decoder 1層 & Spatial Softmax 32 keypoints，U-Netチャンネル(512, 1024, 2048)，DDPM訓練ステップ100 & GPT次元512，8層，head数8，RVQ 2層，codebook size 16 & 状態・行動の最大次元32，flow matching推論10ステップ \\
			学習時メモリ設定 & 32-bit浮動小数点 & 32-bit浮動小数点 & 32-bit浮動小数点 & BF16混合精度，gradient checkpointing有効 \\
			正規化 & 画像・状態・行動: 平均・標準偏差 & 画像: 平均・標準偏差，状態・行動: 最小最大 & 画像: 恒等変換，状態・行動: 最小最大 & 画像: 恒等変換，状態・行動: 平均・標準偏差 \\
			\hline
		\end{tabular}
	\end{table*}
	BF16はbfloat16浮動小数点形式を指す．gradient checkpointingとは，誤差逆伝播時に一部の中間活性を再計算することでGPUメモリ使用量を下げる学習手法である．

	シミュレーションで用いた音響処理パラメータを Table~\ref{tab:audio_params} に示す．
	標準条件では，4アレイ，音情報10FPS，音空間マップとスペクトログラムの両方を用いる．音空間マップに対する追加のガウシアン平滑化，時間平滑化，特徴画像化は行わない．

	\begin{table}[t]
		\centering
		\scriptsize
		\caption{シミュレーション音響処理パラメータ．}
		\label{tab:audio_params}
		\begin{tabular}{lc}
			\hline
			項目 & 値 \\
			\hline
			サンプリング周波数 & 16 kHz \\
			STFT FFT長 / オーバーラップ & 512 / 50\% \\
			音響処理窓 & 1.0 s \\
			マイクロフォンアレイ数 $n$ & 4 \\
			各アレイのマイク数 / 半径 & 8 / 3.5 cm \\
			アレイ配置半径 & 30 cm \\
			部屋サイズ / 反射次数 & $10\,\mathrm{m}\times10\,\mathrm{m}\times3\,\mathrm{m}$ / 3 \\
			MUSIC推定音源数 & 3 \\
			音情報更新頻度 & 10 FPS \\
			NMF成分数 / 閾値 & 50 / $1.6\times10^{-3}$ \\
			スペクトログラム表示帯域 / 正規化 & 0--8000 Hz / min--max \\
			距離減衰 $d_{\mathrm{floor}},\gamma$ & 10.0 m，1.0 \\
			\hline
		\end{tabular}
	\end{table}


	また，Fig.~\ref{fig:appendixA}に，提案手法におけるハイパーパラメータの感度分析結果を示す．

	\begin{figure}[t]
		\centering
		\includegraphics[keepaspectratio, width=14cm]{fig/figb.png}
		\caption{ハイパーパラメータの感度分析．(a) マイクロフォンアレイ数の影響．(b) 音情報の更新FPSの影響．(c) 音空間マップの表現方法の比較．}
		\label{fig:appendixA}
	\end{figure}

\subsection{マイクロフォンアレイ数の影響}
\label{sec:appendixA1}
	SoundPositionタスクに相当する条件において，マイクロフォンアレイ数 $n$ を変化させた際のタスク成功率の変化をFig.~\ref{fig:appendixA}(a)に示す．評価にはACTのOurs3条件を用いた．$n=4$ のとき成功率が最大値を示したため，本論文の実験ではこれを標準設定とした．アレイ数が少ない場合はDOA推定精度が低下し，多すぎる場合は冗長な情報によるモデルの混乱が生じる傾向が見られた．

\subsection{音情報の更新FPSの影響}
\label{sec:appendixA2}
	音空間マップとスペクトログラムの更新頻度を1〜30FPSの範囲で変化させた際の成功率変化をFig.~\ref{fig:appendixA}(b)に示す．評価にはACTのOurs3条件とSoundPositionタスクを用いた．5FPSを超えると性能が急激に向上し，10FPS以降で飽和する傾向が確認された．この結果から，動的な把持動作に対して十分な頻度での音響情報更新が重要であることが示された．

\subsection{音響表現の比較}
\label{sec:appendixA3}
	音空間マップの表現方法として，直接投影，特徴量化マップ，時間平滑音空間マップの3種類を比較した結果をFig.~\ref{fig:appendixA}(c)に示す．評価にはACTとSoundPositionタスクを用いた．直接投影した音空間マップが他の表現を一貫して上回る成功率を達成しており，事前処理を加えることなく生のDOA情報を使用することが有効であることが示された．

\subsection{シミュレーション実験と実機実験の処理差分}
\label{sec:appendixA4}
	実機実験では，シミュレーションと同じMUSIC，音空間マップ生成，Spotforming-NMF，スペクトログラム生成の処理系列を用いる．ただし，実環境のマイクロフォン配置，ロボット構成，実録音に含まれる雑音に合わせて，Table~\ref{tab:sim_real_diff} の処理とパラメータを用いる．
	実機ではSO-100系の単腕構成を用い，右腕7関節のみを状態・行動として扱った．
	そのため，シミュレーションのFranka Pandaにおける9次元状態・9次元行動とは次元が異なる．
	また，実機評価のACTでは時間アンサンブルを無効化し，1回の推論で得たチャンクのうち30ステップを順に実行した．

	\begin{table*}[t]
		\centering
		\scriptsize
		\caption{シミュレーション実験と実機実験の処理差分．}
		\label{tab:sim_real_diff}
		\begin{tabular}{p{0.20\textwidth} p{0.36\textwidth} p{0.36\textwidth}}
			\hline
			項目 & シミュレーション & 実機 \\
			\hline
			ロボット状態・行動 & Franka Pandaの9次元状態・9次元位置制御指令 & 右腕7関節の状態・行動 \\
			カメラ & Genesis内の正面・側面RGB画像を $224\times224$ に保存 & RealSense画像をクロップ・リサイズして $224\times224$ に変換 \\
			音信号 & Pyroomacousticsで音源と部屋反射を生成 & 4台の8chマイクロフォンアレイから16kHzで実録音 \\
			MUSIC推定音源数 & 3 & 1 \\
			DOAスペクトル & MUSIC出力を対数化し，方位方向の総和で正規化 & 相対dB化，60パーセンタイルをノイズ床，15dBをダイナミックレンジとして正規化 \\
			DOA周波数正規化 & なし & 周波数方向に正規化 \\
			音空間マップ座標 & シミュレーション作業空間に対応するローカルグリッド & 実機の矩形マイク配置に合わせた $1.4\,\mathrm{m}$ 四方のマップ座標 \\
			距離減衰 & $d_{\mathrm{floor}}=10.0\,\mathrm{m}$，$\gamma=1.0$ & $d_{\mathrm{floor}}=0.0\,\mathrm{m}$，$\gamma=0.0$（距離減衰なし） \\
			統合マップ & $\alpha=1.0$，局所最大点を使用 & $\alpha=4.0$，最大値位置を使用 \\
			統合マップ平滑化 & ガウシアンフィルタ $\sigma=1.0$ & ガウシアンフィルタ $\sigma=1.0$ \\
			スペクトログラム帯域・正規化 & 0--8000 Hz，min--max正規化 & 0--4000 Hz，2/98パーセンタイル正規化 \\
			音情報更新 & 30FPS環境中で3フレームごとに更新（10FPS） & 1秒窓の音声バッファを10FPSで更新 \\
			チャンネル順序 & シミュレーションのアレイ生成順 & 左上，右上，右下，左下の順序 \\
			評価対象Policy & ACT，Diffusion Policy，VQ-BeT，$\pi_0$ & ACT，Diffusion Policy \\
			実機評価時ACT設定 & シミュレーション評価ではtemporal ensemble係数0.01，実行1 step & temporal ensembleなし，実行30 steps \\
			\hline
		\end{tabular}
	\end{table*}

%===============================================================================
\section{潜在空間の可視化（追加結果）}
\label{sec:appendixB}

	本節では，~\ref{sec:tsne}で示した本文中のFig.~\ref{fig:fig4}に掲載しきれなかったPolicy・タスクの組み合わせについての追加 t-SNE プロットを示す．

\subsection{潜在状態の抽出方法}
\label{sec:appendixB1}
	各Policyは異なる構造を持つため，潜在状態として用いる特徴量も異なる．本研究では，各Policyにおいて行動生成の直前に得られる中間表現を潜在状態として抽出した．具体的には，ACTではTransformer decoderの出力を，Diffusion PolicyではU-Netの最終畳み込み層の直前に位置するボトルネック層直後の特徴量を，VQ-BeTではGPTスタイルのpolicy networkが出力する特徴量のうちaction queryトークンに対応する位置の表現を，$\pi_0$では最終的な行動射影に入力されるsuffix側の出力をそれぞれ用いた．いずれのPolicyも行動チャンクを出力する設計であるため，各エピソードについて行動ステップ方向（チャンク内の時刻方向）に平均をとり，1エピソードあたり1つの潜在ベクトルを得たうえで t-SNE による次元削減を行った．

\subsection{Policy別の追加可視化結果}
\label{sec:appendixB2}
	Fig.~\ref{fig:appendixB_act}〜Fig.~\ref{fig:appendixB_pi0}にACT，Diffusion Policy，VQ-BeT，$\pi_0$それぞれについての全タスクにおける t-SNE プロットを示す．各プロットは，タスク成功・失敗，音源位置，音種をラベルとして色付けしたものである．Ours3条件ではOurs1・Ours2・Baselineに比べて，対応するラベルに応じた構造化された分布が観察されやすく，本文の主要な観察と整合している．

	\begin{figure}[t]
		\centering
		\includegraphics[keepaspectratio, width=14cm]{fig/fig_c_act.png}
		\caption{ACTにおける t-SNE プロット．全4タスク（SoundPosition，SoundTone，SoundAll，SoundShake）について，成功/失敗，音源位置，音種のラベルで色付けした．}
		\label{fig:appendixB_act}
	\end{figure}
	\begin{figure}[t]
		\centering
		\includegraphics[keepaspectratio, width=14cm]{fig/fig_c_diffusion.png}
		\caption{Diffusion Policyにおける t-SNE プロット．}
		\label{fig:appendixB_dp}
	\end{figure}
	\begin{figure}[t]
		\centering
		\includegraphics[keepaspectratio, width=14cm]{fig/fig_c_vqbet.png}
		\caption{VQ-BeTにおける t-SNE プロット．}
		\label{fig:appendixB_vqbet}
	\end{figure}
	\begin{figure}[t]
		\centering
		\includegraphics[keepaspectratio, width=14cm]{fig/fig_c_pi0.png}
		\caption{$\pi_0$における t-SNE プロット．}
		\label{fig:appendixB_pi0}
	\end{figure}

\end{document}
